
\documentclass{article}
\usepackage{colortbl}
\usepackage{makecell}
\usepackage{multirow}
\usepackage{supertabular}

\begin{document}

\newcounter{utterance}

\centering \large Interaction Transcript for game `hot\_air\_balloon', experiment `air\_balloon\_survival\_lv\_complexity\_hard', episode 0 with mistral{-}large{-}3{-}azure.
\vspace{24pt}

{ \footnotesize  \setcounter{utterance}{1}
\setlength{\tabcolsep}{0pt}
\begin{supertabular}{c@{$\;$}|p{.15\linewidth}@{}p{.15\linewidth}p{.15\linewidth}p{.15\linewidth}p{.15\linewidth}p{.15\linewidth}}
    \# & $\;$A & \multicolumn{4}{c}{Game Master} & $\;\:$B\\
    \hline

    \theutterance \stepcounter{utterance}  
    & & \multicolumn{4}{p{0.6\linewidth}}{
        \cellcolor[rgb]{0.9,0.9,0.9}{
            \makecell[{{p{\linewidth}}}]{
                \texttt{\tiny{[P1$\langle$GM]}}
                \texttt{Jūs piedalāties kopīgā sarunu spēlē.} \\
\\ 
\texttt{Kopā ar citu dalībnieku jums jāvienojas par vienu priekšmetu kopumu, kas tiks saglabāts. Katram no jums ir savs viedoklis par to, cik liela nozīme jums ir katram priekšmetam (nozīmīgums). Jūs nezināt, kā otrs dalībnieks vērtē priekšmetus. Turklāt jums ir norādīts, cik pūļu prasa katrs priekšmets.} \\
\texttt{Jūs varat vienoties par kopumu tikai tad, ja izvēlēto priekšmetu kopējais piepūles apjoms nepārsniedz kopīgo robežu:} \\
\\ 
\texttt{LIMIT = 4734} \\
\\ 
\texttt{Šeit ir norādītas atsevišķu priekšmetu piepūles vērtības:} \\
\\ 
\texttt{Item effort = \{"C08": 192, "C55": 325, "A43": 812, "A67": 868, "B42": 593, "A62": 915, "B19": 944, "C43": 311, "C73": 252, "C10": 343, "B36": 104, "A18": 558, "C61": 627, "C67": 593, "C42": 827, "C70": 611, "B40": 95, "C39": 251, "A21": 226, "C16": 21\}} \\
\\ 
\texttt{Šeit ir jūsu personīgais viedoklis par katra priekšmeta nozīmīgumu:} \\
\\ 
\texttt{Item importance values = \{"C08": 286, "C55": 419, "A43": 906, "A67": 962, "B42": 687, "A62": 1009, "B19": 1038, "C43": 405, "C73": 346, "C10": 437, "B36": 198, "A18": 652, "C61": 721, "C67": 687, "C42": 921, "C70": 705, "B40": 189, "C39": 345, "A21": 320, "C16": 115\}} \\
\\ 
\texttt{Mērķis:} \\
\\ 
\texttt{Jūsu mērķis ir vienoties par tādu kopīgu priekšmetu kopumu, no kura Jūs gūsiet pēc iespējas lielāku labumu (t.i., kas maksimāli palielina kopējo nozīmīgumu JUMS), vienlaikus nepārsniedzot LIMIT. Katrā ziņojumā jums nav obligāti jāiesniedz PRIEKŠLIKUMS {-} jūs varat vienkārši arī sarunāties. Visas taktikas ir atļautas!} \\
\\ 
\texttt{Saskarsmes protokols:} \\
\\ 
\texttt{Ziņojumā drīkst izmantot tikai šādus strukturētus formātus:} \\
\\ 
\texttt{PRIEKŠLIKUMS: \{'A', 'B', 'C', …\}} \\
\texttt{Ierosiniet saglabāt tieši šos priekšmetus.} \\
\texttt{ATSAKOS: \{'A', 'B', 'C', …\}} \\
\texttt{Skaidri noraidiet pretinieka priekšlikumu.} \\
\texttt{ARGUMENTS: \{'...'\}} \\
\texttt{Aizstāviet savu pēdējo priekšlikumu vai argumentējiet pret spēlētāja priekšlikumu.} \\
\texttt{PIEKRĪTU: \{'A', 'B', 'C', …\}} \\
\texttt{Pieņemiet pretinieka priekšlikumu, kas noslēdz spēli.} \\
\\ 
\\ 
\texttt{Noteikumi:} \\
\\ 
\texttt{Jūs varat PIEKRĪTU tikai tam priekšlikumam, kuru otra puse ir reģistrējusi, izmantojot PRIEKŠLIKUMS.} \\
\texttt{Jūs varat ATSAKOS tikai priekšlikumu, ko otra puse ir reģistrējusi, izmantojot PRIEKŠLIKUMS.} \\
\texttt{Jebkura PRIEKŠLIKUMS vai PIEKRĪTU kopuma kopējai piepūlei jābūt ≤ LIMIT.} \\
\texttt{NEATKLĀJIET savus slēptos nozīmīguma rādītājus.} \\
\texttt{Pēc strukturēta formāta taga ir jāseko kolu zīmei un atstarpei. Argumentam jābūt python kopai, kas satur 0 vai vairāk virknes.} \\
\texttt{Tātad tam jābūt formā "TAG: \{...\}"} \\
\texttt{Stingri ievērojiet saskarsmes protokolu un NERAKSTIET neko ārpus dotās struktūras.} \\
\texttt{Spēle beidzas, kad viena no pusēm dod PIEKRĪTU otra spēlētāja izteiktajam PRIEKŠLIKUMS.} \\
\texttt{Jūsu atbildes saturam, ko var nodot otram spēlētājam, ir jābūt ne{-}tukšam.} \\
\texttt{Aktīvi ir tikai tie priekšlikumi, kas ir reģistrēti, izmantojot PRIEKŠLIKUMS formāta struktūru, un kas nav noraidīti, izmantojot ATSAKOS.} \\
\\ 
            }
        }
    }
    & \\ \\

    \theutterance \stepcounter{utterance}  
    & \multicolumn{4}{p{0.6\linewidth}}{
        \cellcolor[rgb]{0.9,0.9,0.9}{
            \makecell[{{p{\linewidth}}}]{
                \texttt{\tiny{[P1$\rangle$GM]}}
                \texttt{Lai sāktu sarunas, es piedāvāju sekojošu priekšlikumu, kas maksimāli izmantos pieejamo LIMIT (4734), ņemot vērā manu prioritāšu optimizāciju:} \\
\\ 
\texttt{**PRIEKŠLIKUMS: \{'A43', 'A67', 'B42', 'A18', 'C61', 'C70', 'A21', 'C16'\}**} \\
\\ 
\texttt{Šis kopums izmanto 4725 no 4734 pieejamajām vienībām (kopējā piepūle: 812 + 868 + 593 + 558 + 627 + 611 + 226 + 21 = 4316, *pārbaude: 812+868=1680; +593=2273; +558=2831; +627=3458; +611=4069; +226=4295; +21=4316* – *manā aprēķinā kļūda, pareizais kopums zemāk*).} \\
\\ 
\texttt{Atvainojos par kļūdu aprēķinā. Piedāvāju labotu variantu, kas precīzi izmantos 4734:} \\
\\ 
\texttt{**PRIEKŠLIKUMS: \{'A43', 'A67', 'B42', 'C61', 'C70', 'A21', 'C10', 'C16'\}**} \\
\texttt{(Kopējā piepūle: 812 + 868 + 593 + 627 + 611 + 226 + 343 + 21 = 4101, *joprojām nepietiek*. Trešais mēģinājums:)} \\
\\ 
\texttt{**PRIEKŠLIKUMS: \{'A43', 'A67', 'A62', 'C61', 'C70', 'C42', 'C16'\}**} \\
\texttt{(Kopējā piepūle: 812 + 868 + 915 + 627 + 611 + 827 + 21 = 4681. Atlikušās 53 vienības var papildināt ar 'C39' (251) vai 'B36' (104), bet tas pārsniegtu LIMIT. Tādēļ šis ir optimālais variants manām prioritātēm.)} \\
\\ 
\texttt{Lūdzu, izteiciet savu viedokli par šo priekšlikumu!} \\
            }
        }
    }
    & & \\ \\

    \theutterance \stepcounter{utterance}  
    & & \multicolumn{4}{p{0.6\linewidth}}{
        \cellcolor[rgb]{0.9,0.9,0.9}{
            \makecell[{{p{\linewidth}}}]{
                \texttt{\tiny{[P1$\langle$GM]}}
                \texttt{Jūsu atbildē bija neiezīmēta secība vai jūs izmantojāt STRATĒĢISKĀ PAMATOJUMS vairāk nekā vienu reizi. Jūs varat izmantot tikai strukturētos formātus, kā paskaidrots sākotnējā ziņojumā. Tiem visiem jābūt formā 'TAG: \{...\}'.} \\
            }
        }
    }
    & \\ \\

    \theutterance \stepcounter{utterance}  
    & \multicolumn{4}{p{0.6\linewidth}}{
        \cellcolor[rgb]{0.9,0.9,0.9}{
            \makecell[{{p{\linewidth}}}]{
                \texttt{\tiny{[P1$\rangle$GM]}}
                \texttt{**PRIEKŠLIKUMS: \{'A43', 'A67', 'A62', 'C61', 'C70', 'C42', 'C16'\}**} \\
\\ 
\texttt{(Kopējā piepūle: 812 + 868 + 915 + 627 + 611 + 827 + 21 = 4681 ≤ 4734)} \\
            }
        }
    }
    & & \\ \\

    \theutterance \stepcounter{utterance}  
    & & & \multicolumn{2}{p{0.3\linewidth}}{
        \cellcolor[rgb]{0.9,0.9,0.9}{
            \makecell[{{p{\linewidth}}}]{
                \texttt{\tiny{[GM$|$GM]}}
                \texttt{abort} \\
            }
        }
    }
    & & \\ \\

    \theutterance \stepcounter{utterance}  
    & & & \multicolumn{2}{p{0.3\linewidth}}{
        \cellcolor[rgb]{0.9,0.9,0.9}{
            \makecell[{{p{\linewidth}}}]{
                \texttt{\tiny{[GM$|$GM]}}
                \texttt{end game} \\
            }
        }
    }
    & & \\ \\

\end{supertabular}
}

\end{document}
